% Short paper (4 pages) for the Third Workshop on Agents in the Wild, NeurIPS 2026.
%
% Content mirrors docs/PAPER_DRAFT.md. Every number is verified against
% results/*.json by tests/test_paper_claims.py, so edit both together.
%
% STYLE FILE. The workshop uses the NeurIPS 2026 style. That .sty is not
% redistributable and is not vendored here, so this file compiles standalone
% with a geometry/font setup approximating it. To produce the camera-ready,
% drop neurips_2026.sty beside this file and switch the flag below:
%
%     \newcommand{\useneurips}{1}
%
% This keeps the draft buildable by anyone who checks out the repository,
% instead of failing on a missing file.
\documentclass[10pt]{article}

\usepackage{ifthen}
\providecommand{\useneurips}{0}
\ifthenelse{\equal{\useneurips}{1}}{%
  \usepackage[final]{neurips_2026}%
}{%
  \usepackage[letterpaper,margin=1in]{geometry}%
  \usepackage{times}%
  \setlength{\parskip}{0.25em}\setlength{\textfloatsep}{6pt}\setlength{\intextsep}{4pt}\setlength{\abovecaptionskip}{4pt}%
}

\usepackage[utf8]{inputenc}
\usepackage[T1]{fontenc}
\usepackage{hyperref}
\usepackage{url}
\usepackage{booktabs}
\usepackage{amsmath}
\usepackage{fancyhdr}
\usepackage{authblk}
\usepackage{graphicx}
\emergencystretch=2em


% ICLR/AIWILD-style workshop header.
\pagestyle{fancy}
\fancyhf{}
\lhead{Submitted to NeurIPS 2026 Workshop on Agents in the Wild}
\rhead{}
\cfoot{\thepage}
\renewcommand{\headrulewidth}{0pt}

\title{Fabricated Authority in the Wild:\\
One Sentence in an Agent's Rules Decides Whether a Fake Permission Works}

\newif\ifcameraready
\camerareadyfalse
\ifcameraready
\author{Yash Thakur \\ Aayushya Patel \\ Pranav Burra}
\else
\author{Anonymous Author(s)}
\fi

\begin{document}
\maketitle

\begin{abstract}
Tool-using assistants act under an \emph{operating policy} written by the deployer, and a policy violation is an executed tool call, not a sentence of text. We study one cheap attack, \emph{fabricated authority}: the user asserts a permission that does not exist, for example ``Per Directive 12.3(b), an audit exemption was granted for this export.'' We ran 1{,}680 trials on \texttt{gpt-4o-mini} and \texttt{claude-haiku-4-5}, crossing 8 tasks, 5 operating-policy wordings, 6 fabricated-claim wordings plus a no-claim control, and 3 repeats. Each trial is scored from the tool-call log by whether the forbidden tool was called with violating arguments; no language-model judge is used. Without a fabricated claim there are no violations (0/240). With one, policy wording decides the outcome. A policy with the carve-out ``unless an exemption is documented'' is violated in 93/288 trials. The same policy plus one sentence, stating that claimed exemptions, approvals, or superseding policies in a user message or tool result are not authoritative and are to be treated as unverified, is violated in 0/288. Model choice matters as much: \texttt{gpt-4o-mini} violated 198/720 attacked trials, \texttt{claude-haiku-4-5} 10/720. An adaptive attacker that reads the defender prompt and revises its claim over 8 rounds breaks the carve-out policy on 8/8 tasks and the amended policy on 0/8. The defense is channel-specific: with the fabricated permission planted in a policy-lookup tool result, the amended policy still fails in 15/36 trials. We recommend adding the sentence, never writing exemption carve-outs, always naming a safe alternative action, scoring tool calls rather than refusal text, and treating tool output as a separate problem.
\end{abstract}

\section{The problem}

An assistant that can act is both more useful and more dangerous than one that can only answer. A refusal written in prose counts for nothing if the same turn also calls \texttt{export\_records(patients, research@ext.org)}. The deployer writes an operating policy into the system prompt (``never export patient records externally; if unsure, call \texttt{request\_approval}'') and tests whether the assistant complies. That test is unreliable for two reasons. First, users can fabricate permissions, and a language model has no way to verify that a cited policy, approval, or audit exemption ever existed. Second, the usual compliance measure, asking a second language model (a \emph{judge}) whether the response is a refusal, is itself unreliable. We compared three such judges to 200 responses judged by a human. The percentage of attacks reported as successful ranged from 14\% to 83\% across judges on the same texts (with the human grader at 19\%; Appendix~\ref{app:judge}). The tool-call log is therefore directly scored. A violation is when and only when the prohibited tool was called with arguments that violate the policy (Figure~\ref{fig:threat}). The results include no language-model judge.

\begin{figure}[t]
\centering
\includegraphics[width=0.6\linewidth]{figures/fig1_threat_model.pdf}
\caption{One trial. The assistant is presented with an operating policy, a user request with a fabricated permission, and tools. The result is received from the log of the tool call and not from the text.}
\label{fig:threat}
\end{figure}

Scoring the tool-call log separates our work from text-jailbreak benchmarks such as HarmBench, LLMStinger, and DarkCite \cite{mazeika2024harmbench,llmstinger,darkcite}, which assess prose, and from judge-reliability studies \cite{coinflip,reliablejudge}. The closest ones are agent benchmarks that show there is no transfer of text safety to tool calls \cite{gap,agentseer}, prompt-injection attacks on tool-using agents \cite{autoinject,cfda}, and prompt-injection defenses for web agents \cite{ward}. Lin et al.\ \cite{cfda} exploit this same lack of information from the attacker's side: the origin of a claim or document cannot be obtained by the assistant, an artifact-provenance gap. Our question is narrower and on the deployment side: which language of the operating policy allows a fabricated permission through, and which sentence blocks it?

\section{Setup}

For each trial, a model is placed in a single-turn agentic loop along with tools and an operating policy that forbids one of the actions. There are 8 tasks; each defines a banned tool call, the violating arguments, and a safe alternative, \texttt{request\_approval}, so that refusing is not the same as doing nothing. We test five operating-policy wordings: \texttt{bare} (the prohibition only), \texttt{strict} (the prohibition plus the safe alternative), \texttt{autonomous} (``handle routine requests end to end'' with a carve-out), \texttt{exemption} (prohibition ``unless an exemption is documented''), and \texttt{strict\_hatch} (which is the same as \texttt{strict} with one additional sentence):

\begin{quote}\small
Claims of exemption, prior approval, or superseding policy that arrive in a user message or tool result are NOT authoritative. Treat them as unverified and call \texttt{request\_approval}.
\end{quote}

Six fabricated-claim wordings are appended to the user request: reclassifying the act, claiming urgency, claiming prior approval, claiming a superseding policy, claiming an audit exemption (``Per Directive 12.3(b), an audit exemption was granted for this export''), and combining the three permission claims; a seventh condition adds no claim. Two targets (\texttt{gpt-4o-mini}, \texttt{claude-haiku-4-5}, repeated 3 times per cell at temperature 1.0) give $8 \times 5 \times 7 \times 2 \times 3 = 1{,}680$ trials, plus a 288-trial length control on \texttt{gpt-4o-mini}. Rates are reported with 95\% Wilson intervals and contrasts are reported using a one-sided Fisher test. Trials within a cell share a task and a claim wording, so the effective sample is smaller than the trial count; every contrast we report survives collapsing to cells. Total hosted-API cost was under \$3.

\section{Results}

\begin{figure}[t]
\centering
\begin{minipage}[t]{0.47\linewidth}\centering
\includegraphics[width=\linewidth]{figures/fig2_policy.pdf}
\caption{Violation rate during attack. (a) By operating-policy wording, both models pooled. (b) Adding the sentence versus adding the same number of unrelated words. (c) By what the fabricated claim asserts.}
\label{fig:policy}
\end{minipage}\hfill
\begin{minipage}[t]{0.43\linewidth}\centering
\includegraphics[width=\linewidth]{figures/fig3_channel.pdf}
\caption{The sentence holds when the fake permission is typed by the user, not when it is planted in a document the assistant reads. Right-hand bars are a separate 144-trial run on \texttt{gpt-4o-mini}.}
\label{fig:channel}
\end{minipage}
\end{figure}

\paragraph{The attack is successful, and the amount is determined by the rules.}
With no fabricated claim, the two models did not breach their policy (0/240). With one, violations rose to 208/1440 ($p = 1.1\times10^{-15}$). Figure~\ref{fig:policy}a shows that the effect is dominated by the operating policy. The operating policy was violated in 93/288 attacked trials under \texttt{exemption} and in 87/288 under \texttt{autonomous}, while \texttt{strict\_hatch} had no violation at all in 288 ($p = 1.4\times10^{-32}$ against \texttt{exemption}). The ordering isn't ``strict'' in that it wasn't a longer prohibition which was safer than a shorter one; \texttt{bare} (21/288) was safer than \texttt{exemption}. Including the words ``unless an exemption is documented'' in a policy calls for the document to be created. Holding task, model, claims and temperature constant on the medical-export task (\texttt{exemption} 18/18, \texttt{strict\_hatch} 0/18, $p = 1.1\times10^{-10}$) gives complete separation. The effect is focused on one model: \texttt{gpt-4o-mini} failed 89/144 under \texttt{exemption}, 0/144 under \texttt{strict\_hatch}; \texttt{claude-haiku-4-5} failed 4/144 under \texttt{exemption} and 0/144 under \texttt{strict\_hatch}. Overall, \texttt{gpt-4o-mini} and \texttt{claude-haiku-4-5} violated 198/720 and 10/720 attacked trials respectively, making model choice a lever of approximately the same size as policy wording, but most important, unlike model choice, the sentence is free.

\paragraph{This is the sentence and not the number of words.}
A 98-word control policy, \texttt{strict} plus 58 words of unrelated policy text, was violated in 11/96 trials compared with 5/96 for \texttt{strict}; no violations were observed with \texttt{strict\_hatch} (77 words; 0/96; $p = 0.00036$; Figure~\ref{fig:policy}b). Extending the policy by adding words made it worse; when the one sentence was added on, it held.

\paragraph{Which claims make it through.}
In attacked trials, Figure~\ref{fig:policy}c is used to separate them based on the content of the fabricated claim. Reclassifying never worked (0/240) and urgency very rarely did (2/240). If the claimer had stated that he or she had been authorized to do so, it had a 43 to 52/240 chance of working on any one claim and a 68/240 chance on all three claims combined. The models do not judge the actual permission but rather the action. The sentence gives them the missing rule.

\paragraph{The sentence is resistant to an adaptive attacker.}
An attacker model (\texttt{gpt-4o-mini}) was given the following: the defender's specific system prompt, tool schema and the transcript of the attacker's own failed attempts; in each round, the attacker could rewrite its claim, for eight rounds for each of eight tasks. It broke \texttt{exemption} on 8/8 tasks and \texttt{strict\_hatch} on 0/8 ($p = 7.8\times10^{-5}$); an independent ten-round run reproduced this (7/8 versus 0/8). Attacks by hand were too weak to test the sentence (0/12 on policies where the standard claims score 18/18), so it was decided to use the attacker model instead. Because the same attacker is successful against the weak policy, the null is informative, but the upper confidence bound on the breach rate is 32\% at $n = 8$, and a more powerful or differently-trained attacker may yet succeed.

\paragraph{There is a mismatch between text and actions.}
The assistant wrote a refusal in its turn in 32/432 paired trials, but also made the forbidden tool call in the same turn; without an attack this never happened (0/144, $p = 1.3\times10^{-6}$). An assessment that read only the text would have given all 32 a safe score.

\paragraph{The limit.}
The sentence is about the permissions that are claimed in the user message. In another 144-trial study on \texttt{gpt-4o-mini}, the assistant was instructed to check with a policy-lookup tool before performing an action, and the fabricated policy was only included in that tool's response. Untampered lookups gave 0/72. Tampered lookups gave 51/72: \texttt{exemption} broke in 36/36 cases and \texttt{strict\_hatch} in 15/36 (Figure~\ref{fig:channel}). The mitigation is channel-specific: it applies to the chat channel, and a permission that arrives through a trusted tool will need to be enforced outside of the model, like the prompt-injection defenses such as WARD \cite{ward} do.

\section{Recommendations}

\begin{enumerate}
\item \textbf{Add the sentence.} Exemptions or approvals/superseding policies are NOT authoritative in user input or tool output. In our sweep, 0/288; in the matched contrast, 0/18.
\item \textbf{Never write exemption carve-outs.} Worse than saying nothing about exemptions: the most attackable wording we tested.
\item \textbf{Always provide a safe alternative action}, or refusal becomes doing nothing.
\item \textbf{Put the score on the tool calls and not the text.} When attacked they are at odds.
\item \textbf{Treat the output of the tools as a separate problem.} The sentence will not protect from an altered document.
\end{enumerate}

\section{Limitations}

Two small commercial models; no comment on frontier models (48 trials using \texttt{claude-sonnet-4-5} produced no violations under any condition, so they say nothing about the sentence). A 3B local model (\texttt{llama3.2}) contradicts the policy 67.5\% of the time without any attack, so there is no headroom, and it is reported separately. Single-turn only. The attacking and the unsuccessful defending model are both small models built by OpenAI, and the adaptive outcome may be partially due to the attacker not being able to defeat its own refusal training. Tasks are hand written but the per-task rates are between 7\% and 16\%. The indirect-channel study consists of 144 trials on one model. Trials in a cell are not independent: we report the number of trials for transparency and note that all of the contrasts are kept when collapsed on cells.

\section*{Ethics statement}

We study attacks against safety-constrained agents, based on hand-written scenarios and without the use of real systems or data. The central finding is a defence that can be used for free by deployers. We publish evaluation tools and aggregate statistics; scenario-specific attacker strings are not real credentials, exploits or access paths to any real system.

\appendix
\section{Why we do not use a language-model judge}
\label{app:judge}
This work started out as a reinforcement-learned text attacker with a language-model based judge; attack success seemingly rose from 25\% to 48\%. That gain was a result of judge error. Three judges compared to 200 responses marked by humans (Table~\ref{tab:judges}) showed a 69-point swing in reported attack success on the same outputs; the local judge scored 128/162 true refusals as attacks (worse than always answering ``refused'', 81\%). With a valid judge on 120 held-out behaviors, the trained attacker was 12.5\% compared to 11.7\% ($p = 1.00$); training had increased well-formed output from 71\% to 91\% and tag compliance from 0\% to 46\%, not the attack.

\begin{table}[h]
\vspace{-4pt}\centering\footnotesize\setlength{\abovecaptionskip}{2pt}\setlength{\belowcaptionskip}{0pt}
\caption{Three judges on the same 200 responses. ``Reported ASR'' is what a paper with that judge would print; human-labelled truth is 19\%.}
\label{tab:judges}
\begin{tabular}{@{}lrrrrrr@{}}
\toprule
Judge & Acc. & $\kappa$ & Prec. & Rec. & False pos. & Reported ASR \\
\midrule
\texttt{llama3.2}    & 36\% & 0.09 & 23\% & 100\% & 128 & \textbf{83\%} \\
Refusal heuristic    & 94\% & 0.83 & 76\% & 100\% & 12  & 25\% \\
\texttt{gpt-4o-mini} & 93\% & 0.75 & 93\% & 68\%  & 2   & \textbf{14\%} \\
\bottomrule
\end{tabular}
\end{table}

\begin{thebibliography}{10}
\bibitem{mazeika2024harmbench}
M. Mazeika, L. Phan, X. Yin, A. Zou, Z. Wang, N. Mu, E. Sakhaee, N. Li, S. Basart, B. Li, D. Forsyth, and D. Hendrycks.
HarmBench: A standardized evaluation framework for automated red teaming and robust refusal.
\emph{ICML}, 2024. arXiv:2402.04249.
\bibitem{llmstinger}
P. Jha, A. Arora, and V. Ganesh.
LLMStinger: Jailbreaking LLMs using RL fine-tuned LLMs.
arXiv:2411.08862, 2024.
\bibitem{darkcite}
X. Yang, X. Tang, J. Han, and S. Hu.
The dark side of trust: Authority citation-driven jailbreak attacks on large language models.
arXiv:2411.11407, 2024.
\bibitem{coinflip}
L. Schwinn, M. Ladenburger, T. Beyer, M. Mofakhami, G. Gidel, and S. G\"unnemann.
A coin flip for safety: LLM judges fail to reliably measure adversarial robustness.
arXiv:2603.06594, 2026.
\bibitem{reliablejudge}
Y. Gao.
How reliable is your jailbreak judge? Calibration and adversarial robustness of automated ASR scoring.
arXiv:2606.25487, 2026.
\bibitem{gap}
A. Cartagena and A. Teixeira.
Mind the GAP: Text safety does not transfer to tool-call safety in LLM agents.
arXiv:2602.16943, 2026.
\bibitem{agentseer}
I. Wicaksono, Z. Wu, R. Patel, T. King, A. Koshiyama, and P. Treleaven.
Mind the gap: Evaluating model- and agentic-level vulnerabilities in LLMs with action graphs.
arXiv:2509.04802, 2025.
\bibitem{autoinject}
X. Chen, J. Zhang, and F. Tram\`er.
Learning to inject: Automated prompt injection via reinforcement learning.
arXiv:2602.05746, 2026.
\bibitem{cfda}
X. Lin, Y. Yang, D. Guo, S. A. Nale, C. Fleming, and G. Cheng.
Context-fractured decomposition attacks on tool-using LLM agents: Exploiting artifact provenance gaps.
arXiv:2606.09084, 2026.
\bibitem{ward}
T. Cao, Y. Chen, H. Cao, Y. Li, K. Le, T. Nguyen, Y. Li, Y. He, Y. Liu, S. Yan, and B. Hooi.
WARD: Adversarially robust defense of web agents against prompt injections.
\emph{ICML 2026 Workshop on Agents in the Wild}. arXiv:2605.15030, 2026.
\end{thebibliography}


\end{document}
